% Schémas de Lewis et géométries du méthanal, du cyanure d'hydrogène et du tétrachlorométhane (corrigé du devoir n° 2, exercice 2)
{\renewcommand{\arraystretch}{1}\setlength{\tabcolsep}{2pt}%
\newcommand{\nomL}[1]{{\otLight\fontsize{7.5}{9}\selectfont\color{otMuted}#1}}%
\begin{tabular}{*{3}{>{\centering\arraybackslash}m{44mm}}}
\chemfig{C(-[:150]H)(-[:210]H)=\charge{45=\|,315=\|}{O}} &
\chemfig{H-C~\charge{0=\|}{N}} &
\chemfig{\charge{90=\|,180=\|,270=\|}{Cl}-C(-[2]\charge{0=\|,90=\|,180=\|}{Cl})(-[6]\charge{0=\|,180=\|,270=\|}{Cl})-\charge{0=\|,90=\|,270=\|}{Cl}} \\[1mm]
\nomL{méthanal : plane, angles de \ang{120}} & \nomL{cyanure d'hydrogène : linéaire} & \nomL{tétrachlorométhane : tétraédrique} \\
\end{tabular}}
\par\smallskip\legende{Chaque atome de carbone est entouré de quatre doublets liants ; O, N et Cl complètent leur octet avec des doublets non liants.}
